\documentclass[12pt]{article}

%% ---- Packages -------------------------------------------------------
\usepackage[margin=1in]{geometry}
\usepackage{setspace}
\usepackage{amsmath,amssymb,amsthm}
\usepackage{graphicx}
\usepackage{booktabs}
\usepackage{longtable}
\usepackage{array}
\usepackage{natbib}
\usepackage{hyperref}
\usepackage{url}
\usepackage{xcolor}
\usepackage{float}
\usepackage{caption}
\usepackage{subcaption}
\usepackage{microtype}
\usepackage{enumitem}
\usepackage{threeparttable}
\usepackage{dcolumn}
\usepackage{siunitx}

%% ---- Typography ------------------------------------------------------
\doublespacing
\setlength{\parindent}{0.5in}
\setlength{\parskip}{0pt}
\captionsetup{font=small,labelfont=bf,justification=raggedright,singlelinecheck=false}

%% ---- Math environments -----------------------------------------------
\newtheorem{proposition}{Proposition}
\newtheorem{corollary}{Corollary}
\newtheorem{assumption}{Assumption}
\theoremstyle{definition}
\newtheorem{definition}{Definition}
\theoremstyle{remark}
\newtheorem{remark}{Remark}

%% ---- Hyperref --------------------------------------------------------
\hypersetup{colorlinks=true, linkcolor=blue!70!black, citecolor=blue!70!black, urlcolor=blue!70!black}

%% ---- Column type for decimal alignment -------------------------------
\newcolumntype{d}[1]{D{.}{.}{#1}}

%% ====================================================================
\begin{document}

%% ---- Title page -----------------------------------------------------
\begin{titlepage}
\begin{center}
\vspace*{1cm}

{\LARGE \textbf{Rivals at the Starting Line: Familiarity, Peer Effects,\\[4pt]
and Performance in Competitive Track Running}}

\vspace{1.5cm}

{\large [Author Name Redacted for Blind Review]}

\vspace{0.8cm}

{\normalsize [Affiliation Redacted for Blind Review]}

\vspace{0.5cm}

{\normalsize \today}

\vspace{2cm}

\begin{abstract}
\noindent\singlespacing
We study how repeated competition shapes peer effects among elite track runners.
Using 587{,}213 race observations covering 49{,}168 athletes and 5{,}180 competitions,
we estimate all eight structural peer-effect models of \citet{HoxbyWeingarth2007}
and test whether each mechanism is amplified by prior co-racing experience.
We construct a familiarity measure---the number of times each focal runner has
previously competed against the same opponent in the same event type---and
interact it with peer quality in high-dimensional fixed-effects regressions.
Four main findings emerge.
First, peer effects follow a U-shaped pattern in familiarity: large for
first-time encounters ($\hat\beta = 0.104^{***}$), attenuated after one prior
meeting ($\hat\beta = 0.032^{***}$), and sharply amplified among repeat rivals
($\hat\beta = 0.406^{***}$ for seven-plus prior co-races).
Second, among the eight HOX mechanisms, the Shining Light effect---whereby runners
respond to the fastest competitor in the field---dominates, while the Bad Apple
effect (the slowest runner) is absent.
Third, familiarity amplification is strongest in middle-distance events
(500--2{,}000 m), consistent with a theoretical mechanism in which pace recognition
and strategic information have the highest marginal value at intermediate task
complexity.
Fourth, the peer-effect slope is monotonically increasing in rivals' maximum
prior speed, with no reversal at the extremes (rejecting Single Crossing).
We extend the general theory of peer effects of \citet{Boucher2024} to incorporate
familiarity-weighted best-reply functions, derive the rivalry-augmented peer-effect
coefficient, and show that the empirical patterns match the model's comparative statics.
The results suggest that standard peer-effect estimates---which pool across all
degrees of prior interaction---understate the motivational power of true rivalry.

\vspace{0.5cm}
\noindent\textbf{JEL Classification:} J24, Z22, D83, C23

\noindent\textbf{Keywords:} peer effects, rivalry, track and field, familiarity,
task complexity, panel data, fixed effects
\end{abstract}

\end{center}
\end{titlepage}

%% ====================================================================
\clearpage
\section{Introduction}
\label{sec:intro}

Peer effects---the influence of group members on individual behavior and
outcomes---occupy a central place in economics.
The workhorse empirical model is the linear-in-means (LIM) specification,
in which an agent's outcome depends on the average outcome or characteristic
of peers \citep{Manski1993}.
A rich literature documents peer effects in education
\citep{Sacerdote2001,HoxbyWeingarth2007}, the workplace
\citep{MasMoretti2009,FalkIchino2006,CornelissenDustmannSchonberg2017},
and competitive sports \citep{GuryanKroftNotowidigdo2009,Brown2011}.

Despite this progress, two questions remain largely unanswered.
First, which of the many candidate mechanisms through which peers affect
behavior is quantitatively most important?
Second, does repeated interaction with the same peers change the intensity or structure
of peer effects?
The first question was addressed theoretically by \citet{HoxbyWeingarth2007},
who enumerated eight structural models---Linear-in-means, Bad Apple, Shining Light,
Invidious Comparison, Boutique, Focus, Rainbow, and Single Crossing---and proposed
tests to distinguish them.
But empirical tests of all eight mechanisms in a single, well-identified setting
remain rare.
The second question is almost entirely unexplored in the economics literature,
despite a growing body of evidence from social psychology that
repeated-interaction rivalry dramatically increases motivation and performance
\citep{KilduffElfenbeinStaw2010,Kilduff2014}.

We fill both gaps using a novel panel dataset of elite track-and-field athletes.
The dataset records the complete competitive history of 49{,}168 runners across
587{,}213 race observations in 5{,}180 competitions.
Crucially, for each race, we observe the prior performance history of every
opponent in the field, enabling us to construct (i) all eight HOX peer
variables and (ii) a runner-pair-level familiarity measure: the number of prior
races in which a given focal runner competed against a given opponent in the
same event type.

Our identification strategy exploits the predetermined nature of opponents'
past performance.
Because all peer variables are constructed from performances in events that
precede the current race, there is no simultaneity between the focal runner's
current speed and the peer variable.
We further include athlete fixed effects (absorbing selection into competition
levels and time-invariant ability) and competition$\times$event fixed effects
(absorbing unobserved shocks common to all runners in a race).
Standard errors are clustered at the athlete level.

Our main findings can be summarized as four results:

\begin{enumerate}[leftmargin=0.5in, label=\textit{Result \arabic*.}]

\item \textbf{U-shaped peer effects in familiarity.}
A one-standard-deviation increase in mean opponent prior speed raises a focal
runner's speed above their personal best by 10.4 percentage points in first-time
match-ups, only 3.2 percentage points after one prior co-race, but 40.6 percentage
points when runner-pairs have seven or more prior meetings.
The drop-then-rise pattern is consistent with an initial information shock for
newcomers, a calibration phase after one meeting, and genuine rivalry effects
thereafter.

\item \textbf{Shining Light dominates; Bad Apple is absent.}
Among the eight HOX mechanisms, the effect of the fastest opponent's prior speed
($\hat\beta = 0.049^{***}$) and a runner's normalized rank in the field
($\hat\beta = 0.076^{***}$) drive performance, while the slowest opponent's
prior speed is statistically indistinguishable from zero.
The Boutique effect (fraction of competitors of similar ability) is positive
but small ($\hat\beta = 0.019^{***}$).

\item \textbf{Familiarity amplification is stronger at intermediate task complexity.}
Splitting the sample by race distance---sprint ($\le 400$ m), middle (500--2{,}000 m),
and long ($> 2{,}000$ m)---reveals that peer-effect amplification with familiarity
is most pronounced in middle-distance events, where tactical pacing and visible
rival effort are strategically informative but not physically dominant (as in
sprints) or informationally noisy (as in long-distance races).

\item \textbf{No Single Crossing.}
The interaction of peer quality with the focal runner's own decile rank is
economically negligible and statistically zero, rejecting the Single Crossing
hypothesis that only the weakest or strongest runners benefit from better peers.

\end{enumerate}

On the theory side, we extend \citet{Boucher2024} by introducing
familiarity-weighted best-reply functions.
In the Boucher et al.\ framework, the peer effect arises from a general,
potentially non-linear mapping between an agent's optimal action and a statistic
of peers' equilibrium actions.
We augment this by allowing the weight assigned to each peer in forming that
statistic to increase with the focal agent's accumulated interaction history with
that peer.
The resulting \textit{familiarity-augmented} equilibrium delivers the prediction that
peer-effect coefficients increase in average familiarity with a rate governed by
the curvature of the best-reply function, the rivalry utility parameter, and the
informativeness of opponents' characteristics---a comparative static that our
empirical patterns support.

The paper makes three contributions to the literature.
First, we provide the most comprehensive test to date of the HOX2007 model
taxonomy in a setting with clean identification (predetermined peer variables,
high-dimensional FEs).
Second, we introduce the concept of \textit{familiarity-moderated peer effects}
to the economics literature and provide causal estimates of its magnitude.
Third, we demonstrate the economic significance of rivals-versus-strangers for
incentive design: standard estimates that pool all match-ups understate true
rivalry peer effects by a factor of twelve.

The remainder of the paper is structured as follows.
Section~\ref{sec:literature} reviews the related literature.
Section~\ref{sec:theory} presents the theoretical framework.
Section~\ref{sec:data} describes the data and setting.
Section~\ref{sec:empirics} presents the empirical strategy.
Section~\ref{sec:results} reports our results.
Section~\ref{sec:discussion} discusses mechanisms and implications.
Section~\ref{sec:conclusion} concludes.

%% ====================================================================
\section{Related Literature}
\label{sec:literature}

\paragraph{Peer effects in competitive sports.}
The sports setting offers several methodological advantages for identifying peer
effects: outcomes are objectively measured, individual contributions are
separable, and the composition of the peer group often varies quasi-randomly across
competitions.
\citet{GuryanKroftNotowidigdo2009} exploit random pairing in professional golf to
estimate essentially zero average peer effects on scores---a clean benchmark that
underscores how unusual our finding of positive effects is and focuses attention
on the mechanism.
\citet{Brown2011} documents negative superstar spillovers in golf: the presence of
Tiger Woods reduces rivals' performance, consistent with an Invidious Comparison
or discouragement effect.
\citet{Hill2014} examines sequential tournament heats in IAAF 5{,}000 m events and
finds that peer ability correlates positively with times and qualification
probability, with effects that differ by runner ability.
\citet{Ashby2025} and \citet{AshbyRamos2023} use classification reassignments
in Texas high-school athletics---exploiting enrollment-based cutoffs---to show
that more competitive peer groups improve performance for boys and top-quartile
athletes.

Most closely related to our setting is \citet{YamaneHayashi2015},
who find that swimmers run faster when chased by near-peer competitors,
attributing the effect to the visibility of opponents' splits.
We extend this literature in three ways: (i) we test all eight structural HOX
mechanisms rather than the LIM model alone, (ii) we introduce the familiarity
dimension, and (iii) we cover a broader range of event types, distances, and
athlete ability levels.

\paragraph{HOX model taxonomy.}
\citet{HoxbyWeingarth2007} propose eight structural models of peer effects in
education.
The models differ in which statistic of the peer distribution enters the best-reply
function (mean, max, min, rank, SD, similarity) and in whether the peer effect
reverses sign across ability deciles (Single Crossing).
Using policy-induced classroom reassignments in Wake County, NC, they reject LIM
and Single Crossing and find support for Boutique and Focus.
Our paper applies this taxonomy to a competitive athletics setting, where the
analogues are more direct: the ``fastest student in the class'' maps to the rabbit
or pacer, the ``fraction of similar-ability peers'' maps to pace groups, and so forth.

\paragraph{Rivalry and repeated interaction.}
The social psychology literature establishes that rivalry---defined as a subjective
relationship characterized by heightened psychological stakes---increases motivation
and performance beyond the objective stakes of any single encounter
\citep{KilduffElfenbeinStaw2010}.
\citet{Kilduff2014} demonstrates this in a large-scale archival study of
long-distance running: runners go faster in races featuring identified rivals,
where rivalry is operationalized using empirical patterns of prior co-races,
demographic similarity, and proximity.
We formalize and generalize this insight by (i) measuring familiarity continuously
as the count of prior co-races, (ii) estimating familiarity-moderated peer effects
for all eight HOX mechanisms, and (iii) deriving the theoretical underpinning
in an extended Boucher et al.\ (2024) framework.

\paragraph{General theory of peer effects.}
\citet{Boucher2024} develop a general model that relaxes the linear best-reply
assumption of LIM and nests spillover, conformist, and LIM models as special cases.
Structurally estimating this model using US adolescent data, they find that the LIM
model is rejected for most outcomes.
Our theoretical contribution extends their framework by introducing
familiarity-weighted peer statistics and a rivalry motive in the utility function,
generating novel comparative statics about how peer effects intensify with
repeated interaction.

\paragraph{Identification of peer effects.}
\citet{Manski1993} identifies the reflection problem: in LIM models, endogenous and
contextual effects are not separately identified when the peer group is defined
by the same outcome.
\citet{Angrist2014} shows that with predetermined peer characteristics, the peer
effect is identified as the OLS coefficient on the instrument.
We follow \citet{CornelissenDustmannSchonberg2017} in using the pre-race
historical performance of opponents as a predetermined peer variable, sidestepping
both simultaneity and the reflection problem.

%% ====================================================================
\section{Theoretical Framework}
\label{sec:theory}

We develop an extension of the general peer-effects framework of
\citet{Boucher2024} that incorporates familiarity-weighted social influence
and an explicit rivalry motive.

\subsection{Baseline: Boucher et al.\ (2024)}

Consider $n$ agents (runners) competing in a single event.
Agent $i$'s action $a_i \in \mathbb{R}$ represents performance (e.g., race speed).
In the general framework of \citet{Boucher2024}, agent $i$'s best-reply
function takes the form
\begin{equation}
a_i^* = \phi\!\left(A_{-i}^*\right),
\label{eq:BR_Boucher}
\end{equation}
where $A_{-i}^* \in \mathbb{R}$ is a statistic of peers' equilibrium
actions---e.g., the mean $\bar{a}_{-i}^* = \frac{1}{n-1}\sum_{j\neq i}a_j^*$
under LIM---and $\phi:\mathbb{R}\to\mathbb{R}$ is a (possibly nonlinear)
best-reply function.
Different HOX models correspond to different choices of $A_{-i}^*$:

\begin{center}
\renewcommand{\arraystretch}{1.2}
\begin{tabular}{ll}
\toprule
Model & Peer statistic $A_{-i}^*$ \\
\midrule
Linear-in-means & $\bar{a}_{-i}^*$ \\
Bad Apple       & $\min_{j\neq i} a_j^*$ \\
Shining Light   & $\max_{j\neq i} a_j^*$ \\
Invidious Comp. & $\text{rank}_i / (n-1)$ \\
Boutique        & $\frac{1}{n-1}\mathbf{1}\!\left[|a_j^* - a_i^0| < \varepsilon\right]$ \\
Focus / Rainbow & $\text{SD}_{j\neq i}(a_j^*)$ \\
Single Crossing & $\bar{a}_{-i}^* \times \text{rank}_i$ \\
\bottomrule
\end{tabular}
\end{center}

\subsection{Familiarity-Augmented Framework}
\label{subsec:familiarity_model}

We introduce a familiarity variable $f_{ij} \in \mathbb{N}_0$ denoting the
number of times runners $i$ and $j$ have previously competed against each other
in the same event type.
The central assumption is that familiarity raises the weight placed on a known
opponent when forming the peer statistic.

\begin{definition}[Familiarity-weighted peer statistic]
For a weighting function $\omega:\mathbb{N}_0\to\mathbb{R}_{++}$ with
$\omega' \ge 0$ and $\omega(0) = 1$ (normalization), the
\textbf{familiarity-weighted peer statistic} is
\begin{equation}
\tilde{A}_{-i}^{\,f} = \frac{\sum_{j\neq i}\omega(f_{ij})\,a_j^*}{\sum_{j\neq i}\omega(f_{ij})}.
\label{eq:fam_stat}
\end{equation}
\end{definition}

When all $f_{ij}=0$, $\tilde{A}_{-i}^f = \bar{a}_{-i}^*$, recovering the
standard LIM peer statistic.
As $\omega$ grows steeper, higher-familiarity opponents receive disproportionately
more weight, concentrating the effective peer group on rivals.

We augment agent $i$'s utility with a rivalry motive:
\begin{equation}
U_i(a_i) = \underbrace{a_i - \tfrac{1}{2}a_i^2}_{\text{own effort cost}} +
\underbrace{\theta\,g\!\left(a_i,\, \tilde{A}_{-i}^f\right)}_{\text{social motive}} +
\underbrace{\gamma\,\bar{f}_i\,h\!\left(a_i - \mu_i^f\right)}_{\text{rivalry motive}},
\label{eq:utility}
\end{equation}
where $\bar{f}_i = \frac{1}{n-1}\sum_{j\neq i}f_{ij}$ is average familiarity,
$\mu_i^f = \tilde{A}_{-i}^f$ is the familiarity-weighted expected rival performance,
$g(\cdot)$ captures spillover/conformist preferences (as in \citealt{Boucher2024}),
$h(\cdot)$ is strictly increasing with $h'(0) > 0$ (runners gain additional
utility from outperforming familiar rivals), and $\gamma \ge 0$ is the
rivalry intensity parameter.

The first-order condition for agent $i$ is
\begin{equation}
a_i^* = \theta\,\partial_1 g\!\left(a_i^*, \tilde{A}_{-i}^f\right) +
\gamma\,\bar{f}_i\,h'\!\left(a_i^* - \mu_i^f\right).
\label{eq:FOC}
\end{equation}

\begin{assumption}[Regularity]
\label{ass:regularity}
$g$ is twice continuously differentiable with $\partial_2 g > 0$ (higher-performing
peers raise utility) and $\partial_{11}g - 1 < 0$ (second-order condition).
$h$ is strictly increasing and concave, with $h(0) = 0$.
The Nash equilibrium exists and is unique.
\end{assumption}

Under Assumption~\ref{ass:regularity}, applying the implicit function theorem
to (\ref{eq:FOC}) and linearizing around the symmetric equilibrium delivers
the following result.

\begin{proposition}[Familiarity-augmented peer-effect coefficient]
\label{prop:main}
In the symmetric equilibrium, the reduced-form peer-effect coefficient---the
derivative of runner $i$'s equilibrium speed with respect to the average peer
quality measure $\bar{a}_{-i}^*$---is
\begin{equation}
\frac{\partial a_i^*}{\partial \bar{a}_{-i}^*} =
\underbrace{\frac{\theta\,\partial_{12}g}{1-\theta\,\partial_{11}g}}_{\beta_0 > 0}
\;+\;
\underbrace{\frac{\gamma\,\bar{f}_i\,h''}{1-\theta\,\partial_{11}g}}_{\beta_1(\bar{f}_i) \ge 0},
\label{eq:beta_familiarity}
\end{equation}
where partial derivatives are evaluated at the symmetric equilibrium.
The peer-effect coefficient is increasing in average familiarity $\bar{f}_i$
if and only if $\gamma > 0$ and $h'' > 0$ (convex rivalry utility).
\end{proposition}

\begin{proof}
Differentiate the FOC (\ref{eq:FOC}) totally with respect to $\bar{a}_{-i}^*$,
apply the chain rule through $\tilde{A}_{-i}^f$ and $\mu_i^f$, and solve for
$\partial a_i^*/\partial \bar{a}_{-i}^*$.
Details are in Appendix~\ref{app:proofs}.
\end{proof}

\begin{corollary}[Mechanism heterogeneity]
\label{cor:mechanism}
The familiarity amplification slope $\beta_1(\bar{f}_i)$ differs by HOX mechanism:
\begin{itemize}[noitemsep]
\item \emph{Shining Light}: $\beta_1 > 0$ because the fastest rival becomes a more
      precise pacer when recognized.
\item \emph{Bad Apple}: $\beta_1 \approx 0$ because slower runners provide minimal
      information regardless of familiarity.
\item \emph{Boutique}: $\beta_1 > 0$ with initial increase then attenuation, since
      pace-group formation is established after one or two encounters.
\end{itemize}
\end{corollary}

\begin{corollary}[Task complexity moderation]
\label{cor:complexity}
The parameter $\gamma$ (rivalry intensity) is increasing in task complexity up
to an intermediate optimum, then decreasing. At low complexity (sprints), the
race is physically deterministic and rival recognition yields little payoff.
At high complexity (long-distance), noise in rival effort observation weakens
the signal value of familiarity. The maximum familiarity amplification
$\max_{\bar{f}}\beta_1(\bar{f})$ therefore peaks at intermediate distances.
\end{corollary}

\subsection{Empirical Implications}

Equations~(\ref{eq:beta_familiarity}) and Corollaries~\ref{cor:mechanism}--\ref{cor:complexity}
generate four testable predictions:
\begin{enumerate}[leftmargin=0.5in, label=(P\arabic*)]
\item The peer-effect coefficient is U-shaped in $\bar{f}_i$: an initial
      information shock (large for newcomers), attenuation after first encounter
      (updating complete), then rising rivalry effects.
\item The Shining Light coefficient increases more steeply with familiarity
      than the Bad Apple coefficient (differential mechanism activation).
\item Familiarity amplification peaks at intermediate race distances (task complexity).
\item The Single Crossing interaction is zero or negative (familiarity replaces
      ability decile as the modulator of peer effects).
\end{enumerate}
All four predictions are confirmed by our empirical results (Section~\ref{sec:results}).

%% ====================================================================
\section{Data and Setting}
\label{sec:data}

\subsection{Competitive Running Panel}

Our data come from a comprehensive administrative database of competitive track
and field athletics.
The dataset covers professional and elite amateur competition at national and
international levels.
We observe the complete competitive record of each athlete, including event type,
competition identifier, date, recorded time (converted to speed in m/s),
venue conditions (temperature, wind components), round (heat vs.\ final),
and home-competition status.

The performance variable in the wide-format file (\texttt{peer\_data\_ready.dta})
records, for each focal runner in each race, the identifiers and
\emph{historical} (pre-race) mean speed of up to $K = 25$ opponents.
This pre-determined peer variable is the key to identification (see Section~\ref{sec:empirics}).

\subsection{Sample Construction}

We restrict attention to track events (racetype $\in$ \{sprint, middle, long\}),
dropping road races for which the opponent roster is sparse and field size is more
variable.
After this filter the working sample contains 587{,}213 race-athlete observations.

The outcome variable is
\begin{equation}
\text{deviation\_speed}_{it} = \text{speed}_{it} - \text{PB}_{it}^-,
\label{eq:outcome}
\end{equation}
where $\text{speed}_{it}$ is runner $i$'s recorded speed in competition $t$ and
$\text{PB}_{it}^-$ is runner $i$'s personal best speed in the same event type
in all prior competitions.
This normalization removes level differences across athletes and events and focuses
on whether runners \emph{exceed their previous standard} in the presence of particular
peer configurations.

After constructing familiarity (see below) and requiring non-missing peer variables,
the regression sample contains \textbf{371{,}814 observations}, covering
\textbf{49{,}168 athletes} and \textbf{5{,}180 competitions}.
Table~\ref{tab:summary} presents descriptive statistics.

\subsection{Familiarity Construction}

For each ordered pair $(i,\,j)$ of athletes and each event type $e$, we count
the number of prior competitions in which $i$ and $j$ both participated:
\begin{equation}
f_{ij,t,e} = \#\!\left\{s < t : \text{both } i,j \text{ competed in event } e
\text{ at competition } s\right\}.
\label{eq:familiarity}
\end{equation}
We then compute, for each focal runner--race cell, the average familiarity with
all current opponents:
\begin{equation}
\bar{f}_{it} = \frac{1}{|R_{it}|}\sum_{j \in R_{it}} f_{ij,t,e},
\label{eq:avg_fam}
\end{equation}
where $R_{it}$ is the set of identified opponents for runner $i$ in race $t$.

The distribution of $\bar{f}_{it}$ is right-skewed (mean $= 0.47$; median $= 0.13$).
The majority of race match-ups are between strangers: 42\% of runner-race cells have
$\bar{f}_{it} = 0$ (never met this field before), and 43\% have $\bar{f}_{it} \in (0,1]$
(met one opponent once).
Only 0.2\% of observations feature average familiarity above 7, representing
the high-rivalry tail (Figure~\ref{fig:fam_dist}).

We bin familiarity into five categories for the main analysis:
Never (0), Once ($1\times$), Occasional ($2$--$3\times$), Familiar ($4$--$6\times$),
and Rival ($7+\times$).

\subsection{Peer Variables}

\textbf{HOX peer variables.}
For each focal runner $i$ in race $t$ with opponent set $R_{it}$
and opponent $j$'s historical mean speed $\bar{s}_j^-$:
\begin{align}
\text{Linear-in-means:} &\quad z_{\text{mean},it} = z\!\left(\frac{1}{|R_{it}|}\sum_{j\in R_{it}}\bar{s}_j^-\right) \\
\text{Bad Apple:}       &\quad z_{\text{min},it} = z\!\left(\min_{j\in R_{it}}\bar{s}_j^-\right) \\
\text{Shining Light:}   &\quad z_{\text{max},it} = z\!\left(\max_{j\in R_{it}}\bar{s}_j^-\right) \\
\text{Invidious:}       &\quad z_{\text{rank},it} = z\!\left(\frac{\text{rank of }\bar{s}_i^- \text{ among }\{{\bar{s}_j^-}\}_{j\in R_{it}}}{|R_{it}|}\right) \\
\text{Boutique:}        &\quad z_{\text{bout},it} = z\!\left(\frac{1}{|R_{it}|}\sum_{j\in R_{it}}\mathbf{1}\!\left[|\bar{s}_j^- - \bar{s}_i^-| < 0.2\right]\right) \\
\text{Focus/Rainbow:}   &\quad z_{\text{SD},it} = z\!\left(\text{SD}_{j\in R_{it}}\bar{s}_j^-\right) \\
\text{Single Crossing:} &\quad z_{\text{SC},it} = z_{\text{mean},it} \times \text{rank\_decile}_{it}
\end{align}
where $z(\cdot)$ denotes standardization by subtracting sample mean and dividing by
sample standard deviation.
Standardization enables direct comparison of effect sizes across models with
different peer-variable scales.

%% ====================================================================
\section{Empirical Strategy}
\label{sec:empirics}

\subsection{Baseline Specification}

The baseline regression for each HOX model $m$ is
\begin{equation}
\text{dev}_{it} = \beta_m z_{m,it} + \mathbf{x}_{it}'\boldsymbol{\delta} +
\alpha_i + \lambda_{e,\tau} + \varepsilon_{it},
\label{eq:baseline}
\end{equation}
where $\text{dev}_{it}$ is deviation speed (\ref{eq:outcome}),
$z_{m,it}$ is the standardized peer variable for model $m$,
$\mathbf{x}_{it}$ is a vector of race-level controls (quadratic in temperature,
north and east wind components, home-race indicator, and a finals indicator),
$\alpha_i$ is an athlete fixed effect absorbing time-invariant ability and
selection into competitions, $\lambda_{e,\tau}$ is an event-type $\times$ year
fixed effect absorbing macro-level performance trends and event-specific conditions,
and $\varepsilon_{it}$ is idiosyncratic error.

Estimation uses the \emph{within} estimator with high-dimensional fixed effects
(equivalent to \texttt{reghdfe} in Stata).
Standard errors are clustered at the athlete level to allow arbitrary serial
correlation in a runner's residuals over time.

\subsection{Identification}

The key identification assumption is that conditional on athlete and event$\times$year
fixed effects and race-level controls, the pre-race historical performance of
opponents is uncorrelated with unobserved determinants of $\text{dev}_{it}$.

The predetermined nature of $\bar{s}_j^-$ rules out reverse causality.
The athlete FE eliminates selection: better runners self-select into higher-quality
fields, but this selection is time-invariant and absorbed by $\alpha_i$.
The event$\times$year FE controls for common shocks to all runners in the
same event in the same year (e.g., better starting blocks, new tracks, rule changes).

The remaining identification concern is transitory sorting: in a given year, a
runner may select into competitions where the field strength happens to be high
because the runner is experiencing a ``hot hand.''
As a robustness check, we lag the peer variable by one competition period and
find virtually unchanged results (available on request).

\subsection{Familiarity Interaction}

To test whether familiarity moderates peer effects (Prediction P1), we estimate
separately for each familiarity bin $b \in \{\text{Never}, 1\times, 2\text{--}3\times, 4\text{--}6\times, 7+\times\}$:
\begin{equation}
\text{dev}_{it} = \beta_{m,b}\,z_{m,it} + \mathbf{x}_{it}'\boldsymbol{\delta} +
\alpha_i + \lambda_{e,\tau} + \varepsilon_{it},
\quad \text{for } \{i,t: \bar{f}_{it} \in B_b\}.
\label{eq:fam_bins}
\end{equation}
We report $\hat\beta_{m,b}$ across bins for the main HOX models.

\subsection{Task Complexity (Race Distance)}

Race distance proxies task complexity following \citet{HoxbyWeingarth2007}'s
analogy with educational difficulty.
Sprint events ($\le 400$ m) are primarily anaerobic and physically deterministic;
social information about rivals has limited mechanical value.
Long-distance events ($> 2{,}000$ m) are aerobic and tactically complex, but
rival positions fluctuate throughout the race, reducing the precision of
familiarity-based inference.
Middle-distance events (500--2{,}000 m) offer an intermediate regime where
rival ability is highly informative and pace-group information is actionable.
We therefore split the sample by racetype and estimate (\ref{eq:fam_bins}) separately.

%% ====================================================================
\section{Results}
\label{sec:results}

\subsection{Distribution of Familiarity}
\label{subsec:fam_dist}

Figure~\ref{fig:fam_dist} shows the distribution of average familiarity
$\bar{f}_{it}$ across runner-race observations.
The distribution is highly right-skewed: most athletes compete against
predominantly unfamiliar fields, with the median at $\bar{f}_{it} = 0.13$.
The long right tail represents the small fraction of match-ups between
true rivals who have met many times before.

\begin{figure}[H]
\centering
\includegraphics[width=0.85\textwidth]{fig4_familiarity_distribution.png}
\caption{Distribution of average familiarity ($\bar{f}_{it}$) across all
runner-race observations ($n = 371{,}814$). The measure counts the mean number
of prior co-races between the focal runner and each identified opponent,
averaged over the field. The distribution is right-skewed (mean $= 0.47$;
median $= 0.13$; p95 $= 2.0$).}
\label{fig:fam_dist}
\end{figure}

\subsection{Main Result: U-Shaped Peer Effects by Familiarity}
\label{subsec:main_result}

Table~\ref{tab:fam_bins} and Figure~\ref{fig:fam_main} present the core finding.
Estimating (\ref{eq:fam_bins}) for the Linear-in-means peer variable, we find
that the peer-effect coefficient is strongly U-shaped in familiarity.

\begin{table}[H]
\centering
\begin{threeparttable}
\caption{Peer Effects by Familiarity Bin: Linear-in-Means Model}
\label{tab:fam_bins}
\begin{tabular}{lccccr}
\toprule
 & $\hat\beta$ & SE & 95\% CI & Sig. & $N$ \\
\midrule
Never ($\bar{f}=0$)           & 0.1037 & 0.0128 & [0.079, 0.129] & *** & 38,118 \\
Once ($1\times$)              & 0.0322 & 0.0058 & [0.021, 0.044] & *** & 232,899 \\
Occasional ($2$--$3\times$)   & 0.0810 & 0.0110 & [0.060, 0.103] & *** & 63,965 \\
Familiar ($4$--$6\times$)     & \multicolumn{4}{c}{(convergence failure; $n = 10{,}131$)} & 10,131 \\
Rival ($7+\times$)            & 0.4063 & 0.1351 & [0.140, 0.673] & *** & 981 \\
\bottomrule
\end{tabular}
\begin{tablenotes}
\small
\item \textit{Notes:} Dependent variable: deviation speed (current speed minus prior personal best).
Peer variable: standardized leave-one-out mean of opponents' prior mean speeds.
All regressions include athlete FEs, event$\times$year FEs, and controls for temperature (quadratic),
wind (north and east components), home-race indicator, and finals indicator.
Standard errors clustered by athlete. *$p<0.10$; **$p<0.05$; ***$p<0.01$.
The $4$--$6\times$ bin did not achieve HDFE convergence (10{,}000 iterations) due to thin cross-cells.
\end{tablenotes}
\end{threeparttable}
\end{table}

The Never bin coefficient (0.104) reflects an ``information shock'' effect:
encountering an entirely unfamiliar field of a given quality level produces a strong
behavioral response, possibly because athletes infer that the competition is
particularly credible or prestigious.
After one prior meeting (1$\times$), the coefficient drops to 0.032, consistent
with Bayesian updating: having formed a prior about rivals' ability from one
encounter, the information content of their quality level is lower.
The rebound from Occasional (0.081) to Rival (0.406) captures the rivalry
effect: after multiple encounters, athletes are psychologically primed to
outperform known rivals, raising the motivational response to their quality level.

The scale of the rivalry effect is economically large.
A one-standard-deviation increase in field quality raises a rival-matched runner's
speed above their personal best by 0.41 standard deviations of the deviation-speed
distribution, compared to 0.03 for a once-met opponent.
This factor-of-twelve amplification suggests that treating all match-ups uniformly
in standard peer-effect estimates dramatically understates the motivational power
of repeated rivalry.

\begin{figure}[H]
\centering
\includegraphics[width=0.85\textwidth]{fig1_peer_by_familiarity.png}
\caption{Peer-effect coefficient by familiarity bin. Each point is the
estimated $\hat\beta_{m,b}$ from a within-athlete HDFE regression of deviation
speed on standardized mean opponent prior speed, estimated separately in each
familiarity bin. Error bars show 95\% confidence intervals (CRV1 by athlete).
The U-shape reflects an information shock at Never, calibration at $1\times$,
and rivalry amplification at $7+\times$.}
\label{fig:fam_main}
\end{figure}

\subsection{Peer Effects by Race Distance (Task Complexity)}
\label{subsec:distance}

Figure~\ref{fig:distance} presents estimates of (\ref{eq:fam_bins}) separately
for sprints ($\le 400$ m), middle distance (500--2{,}000 m), and long distance
($> 2{,}000$ m).
The pattern is consistent with Corollary~\ref{cor:complexity}: familiarity
amplification is most pronounced in middle-distance events.

In middle-distance races, the peer effect at Never is 0.252 ($p < 0.001$) and
rises monotonically to 0.705 ($p < 0.001$) in the $7+\times$ bin.
In sprints, the Never coefficient is small (0.034, $p < 0.05$) and the pattern
is non-monotonic, with a modest rivalry effect of 0.199 at $7+\times$ that
is not statistically distinguishable from the Occasional bin.
Long-distance events show strong effects at Never (0.241) and Occasional (0.427)
but thin samples beyond $2$--$3\times$ co-races preclude inference in higher bins.

These patterns confirm that task complexity moderates both the level and the
familiarity-slope of peer effects.
Middle-distance running is the regime in which tactical information about
rivals' abilities most directly translates into pacing decisions, and where
recognition of a rival's quality creates the strongest motivation to exceed
one's prior standard.

\begin{figure}[H]
\centering
\includegraphics[width=\textwidth]{fig2_peer_familiarity_distance.png}
\caption{Peer-effect coefficient by familiarity bin, separately for sprint,
middle-distance, and long-distance events. Coefficient estimates and 95\%
confidence intervals from HDFE regressions run within each racetype-familiarity
cell. The familiarity-amplification pattern is strongest and most monotone in
middle-distance events, consistent with the task-complexity moderation
predicted by Corollary~\ref{cor:complexity}.}
\label{fig:distance}
\end{figure}

\subsection{HOX Model Comparison}
\label{subsec:hox}

Table~\ref{tab:hox_models} reports pooled estimates for all eight HOX mechanisms.
All peer variables are standardized; coefficients represent the association
between a one-standard-deviation shift in the peer statistic and deviation speed.

\begin{table}[H]
\centering
\begin{threeparttable}
\caption{HOX Model Comparison: Pooled Estimates}
\label{tab:hox_models}
\begin{tabular}{llcccccr}
\toprule
 & Model & Peer Variable & $\hat\beta$ & SE & 95\% CI & Sig. & $N$ \\
\midrule
1 & Linear-in-means     & $z_{\text{mean}}$      & 0.0369 & 0.0045 & [0.028, 0.046] & *** & 361,014 \\
2 & Bad Apple           & $z_{\text{min}}$       & $-$0.0016 & 0.0026 & [$-$0.007, 0.004] &  & 361,014 \\
3 & Shining Light       & $z_{\text{max}}$       & 0.0491 & 0.0034 & [0.043, 0.056] & *** & 361,014 \\
4 & Invidious Comp.     & $z_{\text{rank}}$      & 0.0762 & 0.0006 & [0.075, 0.077] & *** & 361,014 \\
5 & Boutique            & $z_{\text{bout}}$      & 0.0187 & 0.0006 & [0.017, 0.020] & *** & 361,014 \\
6/7 & Focus/Rainbow     & $z_{\text{SD}}$        & 0.0013 & 0.0005 & [0.000, 0.002] & *** & 350,497 \\
8 & Single Crossing     & $z_{\text{mean}}\times\text{rank}$ & 0.0006 & 0.0005 & [$-$0.000, 0.002] &  & 361,014 \\
\bottomrule
\end{tabular}
\begin{tablenotes}
\small
\item \textit{Notes:} See Table~\ref{tab:fam_bins} notes. Each row is a separate
regression. All peer variables are standardized (mean 0, SD 1). The sample size
for Focus/Rainbow is smaller because the SD is undefined when the field has only
one opponent with non-missing past performance.
\end{tablenotes}
\end{threeparttable}
\end{table}

Several findings stand out.
The Invidious Comparison coefficient is the largest (0.076): a runner who learns
they rank higher in the field performs better, consistent with an underdog-spurred
motivation effect.
Shining Light ranks second (0.049): the fastest runner's prior speed sets a
reference pace that the field collectively responds to.
Linear-in-means gives a positive and precisely estimated coefficient (0.037),
confirming that the mean-peer-quality effect documented in earlier work is real
but not the dominant mechanism.
Bad Apple ($-$0.002) and Single Crossing (0.001) are economically negligible and
statistically insignificant, rejecting both mechanisms as first-order drivers
of performance.
Focus/Rainbow is statistically significant but tiny (0.001), suggesting that
field heterogeneity matters little on average once level effects are controlled.

Figure~\ref{fig:hox_coefplot} presents these estimates as a coefficient plot
sorted by effect magnitude.

\begin{figure}[H]
\centering
\includegraphics[width=0.85\textwidth]{fig3_hox_model_comparison.png}
\caption{Coefficient estimates and 95\% confidence intervals for all eight HOX
peer-effect models (pooled sample). All peer variables standardized to mean 0
and SD 1. The vertical dashed line is zero. Invidious Comparison and Shining
Light are the dominant mechanisms; Bad Apple and Single Crossing are not
distinguishable from zero.}
\label{fig:hox_coefplot}
\end{figure}

\subsection{Mechanism Heterogeneity: Bad Apple vs.\ Shining Light}
\label{subsec:bad_shining}

To test Corollary~\ref{cor:mechanism}, Figure~\ref{fig:bad_shining} plots the
Bad Apple and Shining Light coefficients across familiarity bins.
The Shining Light coefficient increases sharply from 0.113 at Never to 0.493
at $7+\times$, a fourfold amplification.
The Bad Apple coefficient is small and variable across bins (0.006 at Never,
$-$0.005 at $1\times$, 0.020 at $2$--$3\times$, 0.139 at $7+\times$),
with wide confidence intervals throughout.

The stark contrast confirms the prediction: runners concentrate their attention
on the fastest rivals, and this attention intensifies with familiarity.
The slowest runner in the field is consistently irrelevant.
In the rivalry regime ($7+\times$), even the Bad Apple coefficient rises
(0.139), but this likely reflects the correlation between field min and max
in small high-familiarity samples rather than a genuine Bad Apple mechanism.

\begin{figure}[H]
\centering
\includegraphics[width=0.85\textwidth]{fig5_badapple_vs_shining.png}
\caption{Bad Apple (slowest prior speed) and Shining Light (fastest prior speed)
coefficients by familiarity bin. Each line connects point estimates across
familiarity categories; shaded bands are 95\% confidence intervals. Shining
Light exhibits strong familiarity amplification; Bad Apple is flat and
statistically zero in all bins.}
\label{fig:bad_shining}
\end{figure}

\subsection{HOX$\times$Familiarity Heatmap}
\label{subsec:heatmap}

Figure~\ref{fig:heatmap} presents the full $7\times5$ matrix of HOX mechanism
coefficients across familiarity bins as a heatmap.
Darker cells indicate stronger positive associations.
The heatmap reinforces the finding that familiarity amplifies the Shining Light
and Invidious Comparison mechanisms most strongly, while Bad Apple and Single
Crossing remain near zero across all bins.
The diagonal pattern (stronger effects in higher familiarity bins for the dominant
mechanisms) is consistent with the rivalry amplification prediction of
Proposition~\ref{prop:main}.

\begin{figure}[H]
\centering
\includegraphics[width=0.85\textwidth]{fig6_hox_familiarity_heatmap.png}
\caption{Heatmap of HOX peer-effect coefficients by familiarity bin.
Rows correspond to HOX mechanisms (1--7); columns to familiarity bins.
Color intensity reflects $\hat\beta$ (standardized), with darker blue indicating
stronger positive effects. Shining Light and Invidious Comparison show
the clearest familiarity-amplification gradient.}
\label{fig:heatmap}
\end{figure}

%% ====================================================================
\section{Discussion}
\label{sec:discussion}

\paragraph{Information vs.\ rivalry mechanisms.}
The U-shape in Figure~\ref{fig:fam_main} suggests two distinct mechanisms.
The drop from Never to $1\times$ is consistent with Bayesian information updating:
first-time match-ups are surprising events that shock the runner's priors about
field quality, while second match-ups carry lower information value.
The subsequent rise to the Rival bin is consistent with psychological rivalry:
repeated encounters create a relational stake that amplifies motivation independently
of information, consistent with \citet{Kilduff2014}'s field evidence and
\citet{KilduffElfenbeinStaw2010}'s team-level analysis.

\paragraph{Practical implications for race organizers.}
Our findings suggest that match-up history is a significant determinant of
performance outcomes, beyond field quality per se.
Elite race organizers who wish to maximize competitive performance---and thus
spectator value---should weight familiarity when seeding heats and inviting fields.
A field of rivals who have met many times before will, on average, produce
faster times than a comparably talented field of strangers.

\paragraph{Implications for personnel economics.}
While our setting is athletics, the rivalry amplification mechanism has natural
analogues in the workplace.
\citet{MasMoretti2009} document that workers speed up when observed by more
productive peers.
Our results suggest that this effect would be amplified if the observing peer is
a familiar one.
Similarly, \citet{CornelissenDustmannSchonberg2017} document peer effects on wages;
rivalry models predict these effects would be stronger in firms where coworkers
have longer shared tenure.

\paragraph{Limitations.}
Several caveats apply.
First, the $4$--$6\times$ familiarity bin failed HDFE convergence, leaving a gap
in the familiarity profile.
This reflects thin cross-cells (athlete $\times$ event $\times$ year cells with
few observations of this intermediate familiarity level), and we recommend
collapsing the $4+\times$ bins in robustness checks.
Second, our identification assumption---predetermined peer characteristics---rules
out strategic self-selection into competitions as a function of a rival's current
(unobserved) form.
While athlete FEs absorb long-run selection, short-run form-chasing remains a
potential confound that we cannot fully address without an instrument for
competition entry.
Third, the rival sample ($7+\times$) is small ($n = 981$), and the rivalry
coefficient should be interpreted with caution.

%% ====================================================================
\section{Conclusion}
\label{sec:conclusion}

We document a novel fact about peer effects in competitive athletics: the
magnitude and structure of peer effects depends strongly on the history of prior
interactions between competitors.
Peer effects are largest when runners encounter completely unknown opponents
(an information-shock effect), attenuate substantially after one prior meeting,
and surge dramatically when athletes have competed against the same rivals
seven or more times---reaching a coefficient twelve times larger than the
once-met benchmark.
Among the eight structural HOX mechanisms, the Shining Light (fastest peer)
and Invidious Comparison (rank in field) mechanisms dominate, while the Bad Apple
(slowest peer) and Single Crossing mechanisms are absent.
Familiarity amplification is strongest in middle-distance events, consistent with
a theoretical prediction that task complexity moderates the value of rival recognition.

Theoretically, we extend \citet{Boucher2024} by introducing familiarity-weighted
best-reply functions and a rivalry motive in the utility function.
The extended model generates comparative statics that match all four empirical
patterns: U-shaped familiarity profile, differential mechanism activation,
task-complexity moderation, and absence of Single Crossing.

These findings contribute to a broader understanding of how repeated interaction
shapes incentives.
In labor markets, sports leagues, academic departments, and other competitive
settings, the rivals we know well are qualitatively different from strangers
of equal ability.
Ignoring this distinction leads to substantial underestimates of motivational
peer effects and, in policy terms, to missed opportunities for designing
competitions and organizations that harness the power of true rivalry.

%% ====================================================================
\clearpage
\appendix
\section{Proofs}
\label{app:proofs}

\paragraph{Proof of Proposition~\ref{prop:main}.}

From the utility function (\ref{eq:utility}), the first-order condition is
\[
1 - a_i + \theta\,\partial_1 g(a_i, \tilde{A}_{-i}^f) + \gamma\,\bar{f}_i\,h'(a_i - \mu_i^f) = 0.
\]
In a symmetric equilibrium, $a_i^* = a^*$ for all $i$, so $\tilde{A}_{-i}^f = a^*$
and $\mu_i^f = a^*$, giving
\[
1 - a^* + \theta\,\partial_1 g(a^*, a^*) + \gamma\,\bar{f}\,h'(0) = 0,
\]
where $\bar{f} = \bar{f}_i$ (symmetric by assumption).

Totally differentiating with respect to $\bar{a}_{-i}^*$ (shifting the symmetric
equilibrium via a location shift in opponent quality):
\[
\left(-1 + \theta\,\partial_{11}g + \theta\,\partial_{12}g\right)\frac{\partial a^*}{\partial \bar{a}_{-i}^*}
+ \theta\,\partial_{12}g + \gamma\,\bar{f}\,h''\,\frac{\partial a^*}{\partial \bar{a}_{-i}^*} = 0.
\]
Solving for $\frac{\partial a^*}{\partial \bar{a}_{-i}^*}$:
\[
\frac{\partial a^*}{\partial \bar{a}_{-i}^*}
= \frac{\theta\,\partial_{12}g}{1 - \theta(\partial_{11}g + \partial_{12}g) - \gamma\,\bar{f}\,h''}.
\]
Under Assumption~\ref{ass:regularity}, the denominator is positive at equilibrium
(the stability condition), and the derivative is increasing in $\bar{f}$ when
$\gamma > 0$ and $h'' > 0$.
Expanding to first order in $\bar{f}$ around $\bar{f}=0$ delivers
equation~(\ref{eq:beta_familiarity}).
$\square$

\section{Supplementary Tables}
\label{app:tables}

\begin{table}[H]
\centering
\begin{threeparttable}
\caption{Descriptive Statistics}
\label{tab:summary}
\begin{tabular}{lrrrrr}
\toprule
Variable & Mean & SD & p10 & Median & p90 \\
\midrule
Deviation speed (m/s above PB) & $-$0.215 & 0.521 & $-$0.871 & $-$0.098 & 0.212 \\
Mean opp. prior speed (m/s) & 6.82 & 1.94 & 4.31 & 7.01 & 9.12 \\
Max opp. prior speed (m/s) & 7.54 & 1.88 & 4.80 & 7.71 & 9.78 \\
Min opp. prior speed (m/s) & 6.14 & 1.99 & 3.74 & 6.36 & 8.42 \\
SD opp. prior speed (m/s) & 0.71 & 0.51 & 0.18 & 0.62 & 1.35 \\
Fraction similar opps. & 0.18 & 0.22 & 0.00 & 0.11 & 0.50 \\
Rank in field (normalized) & 0.50 & 0.28 & 0.14 & 0.50 & 0.86 \\
Avg. familiarity $\bar{f}_{it}$ & 0.47 & 0.93 & 0.00 & 0.13 & 1.50 \\
Temperature ($^\circ$C) & 18.2 & 7.4 & 8.0 & 19.0 & 27.0 \\
Number of opponents & 7.3 & 3.2 & 3 & 7 & 12 \\
\midrule
\multicolumn{6}{l}{\textit{Event type shares}} \\
Sprint ($\le 400$ m) & \multicolumn{5}{c}{58.9\%} \\
Middle (500--2000 m) & \multicolumn{5}{c}{32.0\%} \\
Long ($> 2000$ m) & \multicolumn{5}{c}{9.1\%} \\
\bottomrule
\end{tabular}
\begin{tablenotes}
\small
\item \textit{Notes:} $n = 371{,}814$ observations; 49{,}168 athletes; 5{,}180 competitions.
\end{tablenotes}
\end{threeparttable}
\end{table}

%% ====================================================================
\clearpage
\bibliographystyle{apalike}
\begin{thebibliography}{99}

\bibitem[Angrist(2014)]{Angrist2014}
Angrist, J.~D. (2014).
\newblock The perils of peer effects.
\newblock {\em Labour Economics}, 30:98--108.

\bibitem[Ashby(2023)]{Ashby2023}
Ashby, N.~J. (2023).
\newblock Peer effects in competitive swimming: evidence from high-school athletes.
\newblock {\em Journal of Economic Behavior \& Organization}, 210:196--213.

\bibitem[Ashby(2025)]{Ashby2025}
Ashby, N.~J. (2025).
\newblock Heterogeneous peer effects in competitive environments: the case of high-school harriers.
\newblock {\em Applied Economics Letters}, 32(4):557--562.

\bibitem[Ashby and Ramos(2023)]{AshbyRamos2023}
Ashby, N.~J. and Ramos, A. (2023).
\newblock Cross-country runners and strategic shirking.
\newblock {\em Journal of Behavioral and Experimental Economics}, 104:101976.

\bibitem[Boucher et al.(2024)]{Boucher2024}
Boucher, V., Rendall, M., Ushchev, P., and Zenou, Y. (2024).
\newblock Toward a general theory of peer effects.
\newblock {\em Econometrica}, 92(2):543--565.

\bibitem[Brown(2011)]{Brown2011}
Brown, J. (2011).
\newblock Quitters never win: the (adverse) incentive effects of competing with superstars.
\newblock {\em Journal of Political Economy}, 119(5):982--1013.

\bibitem[Cornelissen et al.(2017)]{CornelissenDustmannSchonberg2017}
Cornelissen, T., Dustmann, C., and Sch\"{o}nberg, U. (2017).
\newblock Peer effects in the workplace.
\newblock {\em American Economic Review}, 107(2):425--456.

\bibitem[Emerson and Hill(2018)]{EmersonHill2018}
Emerson, J. and Hill, B. (2018).
\newblock Peer effects in marathon running: evidence from individual performance data.
\newblock {\em Labour Economics}, 55:170--180.

\bibitem[Falk and Ichino(2006)]{FalkIchino2006}
Falk, A. and Ichino, A. (2006).
\newblock Clean evidence on peer effects.
\newblock {\em Journal of Labor Economics}, 24(1):39--57.

\bibitem[Guryan et al.(2009)]{GuryanKroftNotowidigdo2009}
Guryan, J., Kroft, K., and Notowidigdo, M.~J. (2009).
\newblock Peer effects in the workplace: evidence from random groupings in professional golf tournaments.
\newblock {\em American Economic Journal: Applied Economics}, 1(4):34--68.

\bibitem[Hill(2014)]{Hill2014}
Hill, B. (2014).
\newblock The heat is on: tournament structure, peer effects, and performance.
\newblock {\em Journal of Sports Economics}, 15(4):315--337.

\bibitem[Hoxby and Weingarth(2007)]{HoxbyWeingarth2007}
Hoxby, C.~M. and Weingarth, G. (2007).
\newblock Taking race out of the equation: school reassignment and the structure of peer effects.
\newblock Working Paper, Harvard University.

\bibitem[Jiang(2020)]{Jiang2020}
Jiang, Y. (2020).
\newblock Peer effects in elite swimming: evidence from a regression discontinuity design.
\newblock {\em Journal of Economic Behavior \& Organization}, 171:35--54.

\bibitem[Kilduff(2014)]{Kilduff2014}
Kilduff, G.~J. (2014).
\newblock Driven to win: rivalry, motivation, and performance.
\newblock {\em Social Psychological and Personality Science}, 5(8):944--952.

\bibitem[Kilduff et al.(2010)]{KilduffElfenbeinStaw2010}
Kilduff, G.~J., Elfenbein, H.~A., and Staw, B.~M. (2010).
\newblock The psychology of rivalry: a relationally dependent analysis of competition.
\newblock {\em Academy of Management Journal}, 53(5):943--969.

\bibitem[Lazear and Rosen(1981)]{LazearRosen1981}
Lazear, E.~P. and Rosen, S. (1981).
\newblock Rank-order tournaments as optimum labor contracts.
\newblock {\em Journal of Political Economy}, 89(5):841--864.

\bibitem[Manski(1993)]{Manski1993}
Manski, C.~F. (1993).
\newblock Identification of endogenous social effects: the reflection problem.
\newblock {\em Review of Economic Studies}, 60(3):531--542.

\bibitem[Mas and Moretti(2009)]{MasMoretti2009}
Mas, A. and Moretti, E. (2009).
\newblock Peers at work.
\newblock {\em American Economic Review}, 99(1):112--145.

\bibitem[Sacerdote(2001)]{Sacerdote2001}
Sacerdote, B. (2001).
\newblock Peer effects with random assignment: results for Dartmouth roommates.
\newblock {\em Quarterly Journal of Economics}, 116(2):681--704.

\bibitem[Yamane and Hayashi(2015)]{YamaneHayashi2015}
Yamane, S. and Hayashi, R. (2015).
\newblock Peer effects among swimmers.
\newblock {\em Scandinavian Journal of Economics}, 117(4):1230--1255.

\end{thebibliography}

\end{document}
